\chapter{\normalsize METODOLOGI PENELITIAN}

Bab ini berisi uraian mengenai tahapan-tahapan yang ditempuh dalam pelaksanaan penelitian, mulai dari penentuan jenis penelitian, penyusunan rancangan (desain), identifikasi populasi dan sampel, prosedur pengumpulan data, hingga teknik pengumpulan dan analisis data secara komputasional menggunakan algoritma STATIS dan analisis spasial.

%-----------------------------------------------------------------------------%
% --- 3.1 ---
\section{\normalsize Jenis Penelitian}
%-----------------------------------------------------------------------------%
Penelitian ini menggunakan pendekatan \textbf{kuantitatif} dengan jenis penelitian \textbf{dekriptif-eksplanatori}. Secara khusus, penelitian ini memanfaatkan teknik analisis multivariat lanjut (analisis tiga-arah) untuk mengeksplorasi struktur data panel indikator pembangunan manusia di Provinsi Banten, DKI Jakarta, dan Jawa Barat. Pendekatan ini dipilih karena data yang diolah berupa angka-angka statistik yang diperoleh secara sekunder dan akan dianalisis secara objektif menggunakan algoritma statistika untuk menarik kesimpulan mengenai stabilitas dan dinamika wilayah.

%-----------------------------------------------------------------------------%
% --- 3.2 ---
\section{\normalsize Desain Penelitian}
%-----------------------------------------------------------------------------%
Desain penelitian ini dirancang secara sistematis untuk menjawab tantangan data spasio-temporal. Tahapan penelitian dimulai dari input data, pra-pemrosesan, hingga pembentukan hasil akhir berupa visualisasi interaktif. Rancangan jalannya penelitian secara sistematis ditunjukkan pada Gambar \ref{fig:flowchart-penelitian} berikut.

\begin{figure}[H]
\centering
\scalebox{0.82}{%
\begin{tikzpicture}[
    node distance = 0.50cm and 1.6cm,
    startstop/.style  = {ellipse, draw, minimum width=3cm, minimum height=0.8cm,
                          align=center, fill=gray!10, font=\small},
    io/.style         = {trapezium, trapezium left angle=75, trapezium right angle=105,
                          draw, fill=blue!8, align=center, text width=5.5cm, font=\small},
    process/.style    = {rectangle, draw, fill=blue!12, text width=5.5cm,
                          align=center, minimum height=0.9cm, font=\small},
    decision/.style   = {diamond, draw, fill=orange!15, aspect=2.4,
                          text width=3.4cm, align=center, inner sep=1pt, font=\small},
    arrow/.style      = {-Stealth, thick}
]

% ---- Nodes ----
\node (start)  [startstop]                             {Mulai};
\node (input)  [io, below=of start]                    {Input Data\\(Upload dataset, pilih variabel \& tahun)};
\node (eks)    [process, below=of input]               {Eksplorasi Data Deskriptif};
\node (std)    [process, below=of eks]                 {Standardisasi $Z$-score per Tahun};
\node (kmo)    [decision, below=0.9cm of std]          {Uji KMO $\geq 0.5$?};
\node (bart)   [decision, below=1.0cm of kmo]          {Uji Bartlett $p < 0.05$?};
\node (statis) [process, below=1.0cm of bart]          {Analisis STATIS\\$\mathbf{W}_t \;\top\; RV \;\top\; \alpha_t \;\top\; \mathbf{W}_c \;\top\;$ SVD};
\node (perm)   [process, below=of statis]              {Uji Permutasi\\(Signifikansi Stabilitas Struktur)};
\node (sig)    [decision, below=1.0cm of perm]         {Struktur Signifikan?\ $p < 0.05$};
\node (biplot) [process, below=1.0cm of sig]           {Compromise Biplot \& Compromise Loading};
\node (elbow)  [process, below=of biplot]              {Penentuan Jumlah Kluster Optimal\\(Elbow + Silhouette)};
\node (kmeans) [process, below=of elbow]               {K-Means Clustering pada Skor Kompromi};
\node (moran)  [process, below=of kmeans]              {Uji Autokorelasi Spasial\\(Global Moran's~$I$)};
\node (peta)   [process, below=of moran]               {Visualisasi Peta Kluster \& Trajektori Wilayah};
\node (selesai)[startstop, below=of peta]              {Selesai};

% ---- Side nodes ----
\node (ganti)  [process, right=2.0cm of kmo]           {Ganti/Tambah\\Variabel};
\node (tdk)    [startstop, right=2.0cm of sig]         {Struktur Tidak\\Stabil -- Selesai};

% ---- Main flow ----
\draw [arrow] (start)   -- (input);
\draw [arrow] (input)   -- (eks);
\draw [arrow] (eks)     -- (std);
\draw [arrow] (std)     -- (kmo);
\draw [arrow] (kmo)     -- node[left, font=\scriptsize]{Ya}    (bart);
\draw [arrow] (bart)    -- node[left, font=\scriptsize]{Ya}    (statis);
\draw [arrow] (statis)  -- (perm);
\draw [arrow] (perm)    -- (sig);
\draw [arrow] (sig)     -- node[left, font=\scriptsize]{Ya}    (biplot);
\draw [arrow] (biplot)  -- (elbow);
\draw [arrow] (elbow)   -- (kmeans);
\draw [arrow] (kmeans)  -- (moran);
\draw [arrow] (moran)   -- (peta);
\draw [arrow] (peta)    -- (selesai);

% ---- Decision "Tidak" branches ----
\draw [arrow] (kmo)  -- node[above, font=\scriptsize]{Tidak} (ganti);
\draw [arrow] (ganti) |- (input);
\draw [arrow] (bart) -| node[above, font=\scriptsize, pos=0.2]{Tidak} (ganti);
\draw [arrow] (sig)  -- node[above, font=\scriptsize]{Tidak} (tdk);

\end{tikzpicture}%
}% end scalebox
\caption{\textit{Flowchart} Alur Metodologi Penelitian}
\label{fig:flowchart-penelitian}
\end{figure}

%-----------------------------------------------------------------------------%
% --- 3.3 ---
\section{\normalsize Populasi dan Sampel}
%-----------------------------------------------------------------------------%
Populasi dalam penelitian ini adalah seluruh wilayah administratif setingkat kabupaten dan kota di Pulau Jawa. Namun, mengingat ketersediaan data dan fokus analisis pada wilayah metropolitan yang terintegrasi, sampel yang digunakan adalah:
\begin{enumerate}
    \item \textbf{Unit Observasi ($n$):} 41 wilayah yang terdiri dari seluruh kabupaten/kota di Provinsi Banten, DKI Jakarta, dan Jawa Barat.
    \item \textbf{Periode Waktu ($T$):} 5 tahun pengamatan (2020--2024).
    \item \textbf{Teknik Pengambilan Sampel:} Menggunakan \textit{purposive sampling}, di mana wilayah dipilih berdasarkan kriteria kelengkapan data indikator pembangunan manusia yang dipublikasikan secara konsisten oleh BPS selama periode penelitian.
\end{enumerate}

%-----------------------------------------------------------------------------%
% --- 3.4 ---
\section{\normalsize Teknik Pengumpulan Data}
%-----------------------------------------------------------------------------%
Data dalam penelitian ini merupakan data sekunder yang dikumpulkan melalui studi dokumentasi pada publikasi resmi Badan Pusat Statistik (BPS).
\begin{enumerate}
    \item \textbf{Sumber Data:} Website resmi BPS Pusat dan BPS Provinsi (Banten, DKI Jakarta, Jawa Barat) melalui portal Open Data atau publikasi "Provinsi Dalam Angka".
    \item \textbf{Metode Pengumpulan:} Melalui pengunduhan berkas metadata dan tabel statistik dalam format \textit{spreadsheet} (.xlsx dan .csv).
    \item \textbf{Definisi Operasional Variabel:} Agar setiap indikator dapat diinterpretasikan secara konsisten, berikut adalah 11 variabel utama yang digunakan:
    \begin{itemize}
        \item \textbf{IPM:} Indeks Pembangunan Manusia (0--100).
        \item \textbf{PDRB:} Produk Domestik Regional Bruto ADHK 2010 (Juta Rupiah).
        \item \textbf{PPM:} Persentase Penduduk Miskin (\%).
        \item \textbf{GK:} Garis Kemiskinan (Rupiah/Kapita/Bulan).
        \item \textbf{PKKP:} Prevalensi Ketidakcukupan Konsumsi Pangan (\%).
        \item \textbf{TPT:} Tingkat Pengangguran Terbuka (\%).
        \item \textbf{TPAK:} Tingkat Partisipasi Angkatan Kerja (\%).
        \item \textbf{RLS:} Rata-rata Lama Sekolah (Tahun).
        \item \textbf{AHH L/P:} Angka Harapan Hidup (Tahun).
        \item \textbf{IPG:} Indeks Pembangunan Gender (0--100).
        \item \textbf{Gini Ratio:} Indeks ketimpangan pengeluaran (0--1).
    \end{itemize}
\end{enumerate}

%-----------------------------------------------------------------------------%
% --- 3.5 ---
\section{\normalsize Validasi Instrumen}
%-----------------------------------------------------------------------------%
Karena data yang digunakan berasal dari lembaga resmi negara (BPS) yang telah memiliki protokol sensus dan survei (seperti Susenas dan Sakernas) yang terstandardisasi, validitas dan reliabilitas instrumen pengumpulan data primer dianggap sudah terpenuhi. Namun, pada tahap penelitian ini dilakukan \textbf{validasi data sekunder} melalui:
\begin{enumerate}
    \item \textbf{Penyelarasan Metadata:} Memastikan definisi operasional variabel tetap konsisten di seluruh provinsi dan tahun.
    \item \textbf{Pembersihan Data (\textit{Data Cleaning}):} Melakukan pemeriksaan terhadap anomali atau nilai yang hilang (\textit{missing values}). Wilayah atau variabel yang memiliki data tidak lengkap dikeluarkan untuk menjaga integritas matriks tiga-arah.
    \item \textbf{Standardisasi $Z$-score:} Untuk memvalidasi bahwa setiap variabel dapat diperbandingkan tanpa dipengaruhi oleh perbedaan satuan (skala), diterapkan transformasi:
    \begin{equation}
        z_{ijt} = \frac{x_{ijt} - \bar{x}_{jt}}{s_{jt}}
    \end{equation}
\end{enumerate}

Untuk memastikan data layak dianalisis secara multivariat, dilakukan pengujian berikut:
\begin{enumerate}
    \item \textbf{Indeks Kaiser-Meyer-Olkin (KMO):} Mengukur tingkat kecukupan sampel dengan rumus:
    \begin{equation}
        KMO = \frac{\sum \sum_{j \neq k} r_{jk}^2}{\sum \sum_{j \neq k} r_{jk}^2 + \sum \sum_{j \neq k} a_{jk}^2}
    \end{equation}
    dengan $r_{jk}$ adalah koefisien korelasi sederhana dan $a_{jk}$ adalah koefisien korelasi parsial. Nilai $KMO > 0,5$ menunjukkan data layak dianalisis.
    
    \item \textbf{Uji Bartlett’s Test of Sphericity:} Menguji apakah matriks korelasi merupakan matriks identitas melalui statistik uji:
    \begin{equation}
        \chi^2 = -\left( n - 1 - \frac{2p + 5}{6} \right) \ln |R|
    \end{equation}
    dengan $n$ adalah jumlah sampel, $p$ jumlah variabel, dan $|R|$ adalah determinan matriks korelasi. Jika $p\text{-value} < 0,05$, maka terdapat korelasi antar variabel yang signifikan.
\end{enumerate}

%-----------------------------------------------------------------------------%
% --- 3.6 ---
\section{\normalsize Teknik Analisis Data}
%-----------------------------------------------------------------------------%
Teknik analisis data mengacu pada algoritma STATIS (\textit{Structuration des Tableaux à Trois Indices de la Statistique}) untuk mengolah data panel serta analisis autokorelasi spasial. 

\subsection{\normalsize Ilustrasi Berbasis Data Sampel}
Sebelum menguraikan seluruh langkah algoritma secara teoritis, berikut disajikan ilustrasi perhitungan menggunakan data sampel dari file \texttt{data2\_sample.xlsx} yang mencakup 4 wilayah dan 3 variabel pada tahun 2020. Ilustrasi ini bertujuan untuk memperjelas jalannya rumus-rumus statistika yang digunakan.

\textbf{1. Data Mentah (Tahun 2020):}
Misalkan kita memiliki 4 wilayah: (A) Jakarta Pusat, (B) Bandung Barat, (C) Tasikmalaya, dan (D) Cilegon dengan 3 indikator (PPM, RLS, TPT).
\[
\mathbf{X}_{2020} = 
\begin{pmatrix}
4,51 & 11,38 & 10,97 \\
10,49 & 8,19 & 12,25 \\
12,97 & 9,33 & 7,99 \\
3,69 & 9,87 & 12,69
\end{pmatrix}
\]

\textbf{2. Standardisasi $Z$-score:}
Setelah dihitung rata-rata dan simpangan bakunya, diperoleh matriks standar $\mathbf{Z}_{2020}$ sebagai berikut:
\[
\mathbf{Z}_{2020} = 
\begin{pmatrix}
-0.87 & 1.47 & 0.00 \\
0.66 & -1.31 & 0.69 \\
1.29 & -0.32 & -1.63 \\
-1.08 & 0.15 & 0.93
\end{pmatrix}
\]

\textbf{3. Matriks Intra-struktur ($\mathbf{W}_t$):}
Matriks kemiripan diperoleh dari $\mathbf{W} = \mathbf{Z}\mathbf{Z}^T$ yang menunjukkan kaitan antar-wilayah:
\[
\mathbf{W}_{2020} = 
\begin{pmatrix}
2,91 & -2,50 & -1,58 & 1,16 \\
-2,50 & 2,63 & 0,13 & -0,26 \\
-1,58 & 0,13 & 4,40 & -2,95 \\
1,16 & -0,26 & -2,95 & 2,06
\end{pmatrix}
\]

\textbf{4. Matriks Inter-struktur ($R$) dan Bobot ($\alpha$):}
Dengan membandingkan matriks $\mathbf{W}$ untuk tahun 2020, 2021, dan 2022 menggunakan koefisien RV, diperoleh matriks korelasi vektoriel:
\[
\mathbf{R} = 
\begin{pmatrix}
1,000 & 0,894 & 0,872 \\
0,894 & 1,000 & 0,993 \\
0,872 & 0,993 & 1,000
\end{pmatrix}
\quad \to \quad
\alpha = 
\begin{pmatrix}
0,561 \\ 0,587 \\ 0,583
\end{pmatrix}
\]
Bobot $\alpha$ terbesar menunjukkan tahun yang paling representatif terhadap struktur rata-rata (kompromi).

\textbf{5. Matriks Skor Kompromi ($\mathbf{F}_c$):}
Melalui dekomposisi nilai singular (SVD) pada matriks kompromi $\mathbf{W}_c = \sum \alpha_t \mathbf{W}_t$, diperoleh koordinat wilayah pada ruang dua dimensi:
\[
\mathbf{F}_c = 
\begin{pmatrix}
-2,41 & 0,37 \\
2,39 & -1,01 \\
0,96 & 2,04 \\
-0,95 & -1,41
\end{pmatrix}
\]
Koordinat inilah yang dipetakan sebagai titik-titik fundamental dalam biplot.

\subsection{\normalsize Algoritma STATIS Lengkap}
Setelah memahami ilustrasi di atas, prosedur sistematis yang diterapkan pada dataset utama adalah:
\begin{enumerate}
    \item \textbf{Analisis Inter-struktur (Koefisien RV):} Mengukur kemiripan struktur antar-tahun $t$ dan $t'$ menggunakan rumus:
    \begin{equation}
        RV(\mathbf{W}_t, \mathbf{W}_{t'}) = \frac{\text{tr}(\mathbf{W}_t^T \mathbf{W}_{t'})}{\sqrt{\text{tr}(\mathbf{W}_t^T \mathbf{W}_t) \text{tr}(\mathbf{W}_{t'}^T \mathbf{W}_{t'})}}
    \end{equation}
    Keanggotaan matriks RV membentuk matriks korelasi $\mathbf{R}$ yang kemudian dicari nilai eigen utamanya ($\lambda$) melalui persamaan:
    \begin{equation}
        \mathbf{R}\boldsymbol{\alpha} = \lambda \boldsymbol{\alpha}
    \end{equation}
    dengan $\boldsymbol{\alpha} = (\alpha_1, \dots, \alpha_T)^T$ adalah vektor bobot optimal yang dinormalisasi ($\sum \alpha_t^2 = 1$).

    \item \textbf{Dekomposisi Struktur Kompromi (SVD):} Matriks kompromi $\mathbf{W}_c = \sum \alpha_t \mathbf{W}_t$ didekomposisi secara spektral:
    \begin{equation}
        \mathbf{W}_c = \mathbf{V} \mathbf{\Lambda} \mathbf{V}^T
    \end{equation}
    Skor koordinat wilayah (skor kompromi) diperoleh melalui $\mathbf{F}_c = \mathbf{V} \mathbf{\Lambda}^{1/2}$.

    \item \textbf{Uji Stabilitas (Permutasi):} Kekuatan struktur kompromi divalidasi dengan membandingkan nilai $\lambda$ asli terhadap distribusi $\lambda$ hasil pengacakan (\textit{permutation test}) sebanyak 1000 kali.
\end{enumerate}

\subsection{\normalsize Analisis Autokorelasi Spasial (Global Moran's~\texorpdfstring{$I$}{I})}
Sejalan dengan hasil pengelompokan (klastering K-Means) pada skor kompromi, dilakukan pengujian untuk mendeteksi apakah terdapat ketergantungan spasial menggunakan rumus:
\begin{equation}
    I = \frac{N}{\sum_i \sum_j w_{ij}} \frac{\sum_i \sum_j w_{ij} (x_i - \bar{x})(x_j - \bar{x})}{\sum_i (x_i - \bar{x})^2}
\end{equation}
dengan $w_{ij}$ adalah matriks pembobot spasial berbasis \textit{Queen Contiguity}. Signifikansi statistik $I$ diuji menggunakan skor-$Z$ dengan bantuan nilai harapan $E(I)$ dan varians $Var(I)$:
\begin{equation}
    Z(I) = \frac{I - E(I)}{\sqrt{Var(I)}}, \quad E(I) = \frac{-1}{n-1}
\end{equation}

\subsection{\normalsize Pengembangan Visualisasi Interaktif}
Algoritma di atas dibungkus dalam sistem \textit{dashboard} menggunakan Python (Dash/Plotly) yang memungkinkan:
\begin{itemize}
    \item Proyeksi dinamis (\textit{trajectory}) posisi wilayah setiap tahun ke ruang kompromi.
    \item Rotasi sumbu (\textit{Baseline Rotation}) untuk mempermudah interpretasi ekonomi.
\end{itemize}
